\documentclass[a4paper,12pt]{article}
\usepackage{bbding,combelow,textcomp,amsfonts,amsthm,amsmath,amssymb,graphicx,color,hyperref,etoolbox}
\usepackage[utf8x]{inputenc}
\usepackage[lined,boxed,commentsnumbered]{algorithm2e}

\newtoggle{story}


%%%%%%%%%%%%%%%%%%%%%%%%%%%%%%%%%%%%%%%%%%
%% Customisation options --- update as necessary
%%%%%%%%%%%%%%%%%%%%%%%%%%%%%%%%%%%%%%%%%%

%% Course details: title, semester, instructors

\newcommand{\courseTitle}{Combinatorics II (Math 7702)}
\newcommand{\semester}{Spring 2026}
\newcommand{\instructorA}{Shagnik Das}
\newcommand{\instructorB}{Lu Yun-Chi}

%% Sheet number

\newcommand{\sheetNumber}{3}


%% Submission details

\newcommand{\dueDate}{15:30, May 5th.}
\newcommand{\lateWarning}{Late submissions will receive no credit.}


%% Display irrelevant footnotes?  \toggletrue for yes, \togglefalse for no

\toggletrue{story}

%%%%%%%%%%%%%%%%%%%%%%%%%%%%%%%%%%%%%%%%%%
%% End of customisation options
%%%%%%%%%%%%%%%%%%%%%%%%%%%%%%%%%%%%%%%%%%

\pdfpagewidth 8.5in
\pdfpageheight 11in
\topmargin -1in
\headheight 0in
\headsep 0in
\textheight 8.5in
\textwidth 6.5in
\oddsidemargin 0in
\evensidemargin 0in
\headheight 77pt
\headsep 0in
\footskip .75in

\makeatletter
\newcommand{\buildtitle}[4]{
\begin{flushleft}
{\large
#1
\hfill{}
#2
\par
#3
}
\end{flushleft}
\vskip 4pt
\begin{center}
{\large\bfseries#4\par}
\end{center}
\bigskip
}
\makeatother

\renewcommand{\thesection}{\normalsize\arabic{section}}

\renewcommand{\deg}{\textrm{deg}}
\newcommand{\eps}{\varepsilon}
\newcommand{\ceil}[1]{\left \lceil #1 \right \rceil}

\newcommand{\hint}{\begin{flushright} [Hint at \url{\hintURL}.] \end{flushright}}

\newcommand{\card}[1]{\left| #1 \right|}
\newcommand{\mb}[1]{\mathbb{#1}}
\newcommand{\mc}[1]{\mathcal{#1}}


\newcommand{\story}[1]{\iftoggle{story}{\footnote{#1}}{}}
\newcommand{\storymark}[1][42]{\iftoggle{story}{\footnotemark[#1]}{}}
\newcommand{\storytext}[2][42]{\iftoggle{story}{\footnotetext[#1]{#2}}{}}

\newcommand{\bonus}[2]{\paragraph{Bonus (#1 pt\ifstrequal{#1}{1}{}{s})} #2}


\begin{document}


\buildtitle{\courseTitle{}}{\semester}{\instructorA{} \hfill \instructorB{}}{Exercise Sheet \sheetNumber{}}

\vspace{-0.2in}

\begin{center}

{B13902024 Chang, Yi-Kuei \\ \bf Due date: \dueDate \\ \lateWarning}

\end{center}

\noindent You should try to solve all of the exercises below, and then submit two solutions to be graded --- each problem is worth $10$ points. Make sure to clearly justify your answers, defining everything precisely and stating any results you are using. You should then submit your solutions on NTU COOL, either as a typed PDF or a clear picture or scan of your handwritten answers.

\paragraph{Exercise 1}  Show that the stable matching from the Gale--Shapley algorithm is the best-case outcome for all the men and the worst-case outcome for all the women. More precisely, show that there is no stable matching that pairs any man with a better partner than in the Gale--Shapley matching, and there is no stable matching that pairs any woman with a worse partner than in the Gale--Shapley matching.

\paragraph{Solution:} Suppose the matching generated by Gale-Shapley is \(S^*\). In the following text, we say \(A\) is a valid partner of \(B\) if there is a stable matching that matches \(A\) and \(B\). 

\textbf{claim:} \(S^*\) is the best-case outcome for all the men. 

\textbf{proof:} Suppose not, then there is a man in \(S^*\) who is matched with someone other than his best valid partner. Hence, there is a man who was rejected by his best valid partner in Gale-Shapley since each man move down his preference list in Gale-Shapley. In particular, there is a man who was rejected by a valid partner in Gale-Shapley. Let \(Y\) be the first man rejected by one of his valid partner in Gale-Shapley, and \(A\) be the first valid partner of \(Y\) rejecting him. Thus, when \(Y\) is rejected by \(A\), there is another man \(Z\) such that \(A\) says "maybe" to him. Note that \(A\) prefers \(Z\) to \(Y\). Now let \(S\) be a stable matching where \(A\) and \(Y\) are matched. Let \(B\) be the partner of \(Z\) in \(S\). Since \(Y\) is the first man who gets rejected by a valid partner, \(Z\) is not rejected by any of his valid partner when \(Z\) proposes to \(A\). Note that \(B\) is a valid partner of \(Z\), so \(Z\) hasn't proposed to \(B\) when \(Z\) proposed to \(A\). This shows \(Z\) prefers \(A\) to \(B\) since each man moves down his preference list in Gale-Shapley. Hence, \(A\) and \(Z\) form a unstable pair in \(S\), contradicting the stability of \(S\).

\textbf{claim:} \(S^*\) is the worst-case outcome for all the women. 

\textbf{proof:} Suppose not, then for some woman \(A\) who is matched with a man \(Y\) in \(S^*\), we know there is a stable matching \(S\) s.t. \(A\) is matched with another man \(Z\) in \(S\), where \(A\) prefers \(Y\) to \(Z\). Suppose \(Y\) is matched with \(B\) in \(S\). Since we have shown that \(S^*\) is best-case outcome for all men, so \(Y\) prefers \(A\) to \(B\). This shows \(A\) and \(Y\) form an unstable pair in \(S\), contradicting the stability of \(S\).            

\paragraph{Exercise 2}  
\vspace{-0.2cm}
\begin{itemize}
	\item[(a)] Prove that for any graph $G$, $\kappa(G) \le \kappa'(G)$.
	\item[(b)] For any $1 \le k \le \ell \le m$, construct a graph $G$ with $\kappa(G) = k$, $\kappa'(G) = \ell$ and $\delta(G) = m$.
\end{itemize}

\paragraph{Solution:}
\begin{itemize}
	\item [(a)] Let \(F \subseteq E(G)\) where \(\vert F \vert = \kappa ^{\prime} (G) \) s.t. \(G - F\) is disconnected. Let \(F = [S, \overline{S}]\), where \(\varnothing \neq S \subsetneq V\) and \(\overline{S} = V \setminus S\). If \(\exists x \in S\) and \(y \in \overline{S}\) s.t. \(\left\{ x, y \right\} \notin E(G) \), then for every \(\left\{ u, v \right\} \in [S, \overline{S}] \), if \(u \neq x\), remove \(u\), otherwise remove \(v\). This ensures \(x, y\) are in the remaining graph.  Hence, we can remove \(\le \vert F \vert \) vertices to make remove all edges in \(F\), and at the end, \(x, y\) are disconnected. Hence, \(\kappa (G) \le \vert F \vert = \kappa ^{\prime} (G) \). Otherwise, if \(\left\{ x, y \right\} \in E(G) \) for all \(x \in S\) and \(y \in \overline{S}\), then
    \[
        \vert F \vert = \vert S \vert \cdot \vert \overline{S} \vert = \vert S \vert \cdot (n - \vert S \vert ) \ge n - 1. 
    \]                
    Hence, \(\kappa (G) \le n - 1 \le \vert F \vert = \kappa ^{\prime} (G) \).
	\item [(b)] Suppose initially our graph \(G\) have two parts, say \(H_1, H_2\), where \(H_1, H_2\) are two distinct \(K_{m+1}\) and no edges between them. Now we pick \(k\) vertices \(\left\{ u_1, u_2, \dots , u_k \right\} \subseteq V(H_1) \) and another \(k\) vertices \(\left\{ w_1, w_2, \dots , w_k \right\} \subseteq V(H_2) \), then add \(k\) edges \(\left\{ u_i, w_i \right\}  \) for all \(1 \le i \le k\) to \(G\). Next, we connect some edges from \(\left\{ u_1, \dots , u_k \right\} \) to \(H_2\) so that there are \(\ell\) edges (including \(\left\{ u_i, w_i  \right\} \)'s) between \(H_1\) and \(H_2\). Note that we can always do this since \(\ell \le m\). Now in \(G\) we know \(\delta (G) = m\), and \(\kappa (G) = k\) since \(\left\{ u_1, u_2, \dots , u_k \right\} \) is a separating set of vertcies. Besides, \(\kappa ^{\prime} (G) = \ell \) since we need to remove all \(\ell \) edges between \(H_1\) and \(H_2\) to disconnect \(G\).                
\end{itemize}

\paragraph{Exercise 3}  Prove that if $\Delta(G) \le 3$, then $\kappa(G) = \kappa'(G)$.

\paragraph{Exercise 4}   Let $G$ be a graph and let $u, v$ be distinct vertices in $G$. Recall that, if $u$ and $v$ are nonadjacent, then $\kappa(u,v)$ is the size of the smallest set of vertices whose removal disconnects $u$ from $v$, while $\lambda(u,v)$ is the maximum size of a set of pairwise internally-disjoint $u,v$-paths.

Similarly (even when $u$ and $v$ are adjacent), we define $\kappa'(u,v)$ to be the smallest set of \emph{edges} whose removal disconnects $u$ from $v$, and $\lambda'(u,v)$ to be the maximum size of a set of pairwise edge-disjoint $u,v$-paths.

By constructing an auxiliary graph $H$ whose vertices are the edges of $G$, and using the local vertex version of Menger's Theorem, prove that $\kappa'(u,v) = \lambda'(u,v)$.


% \newpage

% \section*{Hints}  The following QR codes contain hints for some of the homework exercises.  You should be able to decode them using any QR code scanner, including this one: \url{https://zxing.org/w/decode.jspx}.

% \paragraph{Exercise ?} Also available at \url{https://i.ibb.co/???????/S??E?.png}.

\begin{center}
% \includegraphics[scale=0.5]{Hints/S??E?.png}
\end{center}

\end{document}